\section{極形式}
  \subsection{極座標系の導入}
    前節で確認した三角関数を用いると,平面上の点を距離と角度によって表すことができる.
    普段用いている直交座標系と比較しながら,この表し方を導入する.
    直交座標系とは,デカルト座標系とも呼ばれ,平面上で直交する2直線の交点を基準として任意の点を表す手法のことである.
    しかしながら円や媒介変数表示された関数などにおいて,一部煩雑な表現をせざるをえない状況があった.
    また,図形の平行移動や対称移動に関しては容易に行えたものの,回転させたり歪ませたりするような操作は非常に複雑な操作が必要だった.
    そこで,直交座標系に加え,新たな概念として極座標系の導入を行う.
    
    極座標系とは,平面上の任意の点を原点からの距離と角度を用いて記述する方法である.
    例えば

    \begin{tikzpicture}[scale=0.8]
      % 座標軸
      \draw[->,>=stealth,semithick] (-1.5,0)--(5.5,0) node[above]{$x$};
      \draw[->,>=stealth,semithick] (0,-1)--(0,5) node[right]{$y$};

      % 原点
      \draw (0,0) node[below right]{O};

      % 点 (1,sqrt(3)) への補助線
      \draw[dashed] (1,0) -- (1,{sqrt(3)});
      \draw[dashed] (0,{sqrt(3)}) -- (1,{sqrt(3)});

      % 座標軸上の値
      \draw (1,0) node[below]{$1$};
      \draw (0,{sqrt(3)}) node[left]{$\sqrt{3}$};

      % 点
      \fill (1,{sqrt(3)}) circle (2pt)
        node[above right]{$(1,\sqrt{3})$};
    \end{tikzpicture}

    のように表される点は,極座標系においては

\begin{tikzpicture}[scale=0.8]

  % 座標軸
  \draw[->,>=stealth,semithick]
    (-1.5,0)--(5.5,0) node[above]{$x$};
  \draw[->,>=stealth,semithick]
    (0,-1)--(0,5) node[right]{$y$};

  % 原点・点A
  \coordinate (O) at (0,0);
  \coordinate (A) at (1,{sqrt(3)});
  \node[below left] at (O) {O};

  % 動径
  \draw[very thick] (O)--(A);

  % 下側の曲線
  \draw[semithick]
    (O)
    .. controls (0.02,0.35) and (0.12,0.62)
    .. (0.3,0.7);

  % 上側の曲線
  % 終点Aでの接線を動径OAと平行にする
  \draw[semithick]
    (0.6,1.4)
    .. controls (0.62,1.58) and
               (0.82,{sqrt(3)-0.10*sqrt(3)})
    .. (A);

  % 動径の長さ
  \node at (0.30,1.05) {$2$};

  % 角度
  \draw[semithick]
    (0.55,0) arc (0:60:0.55);
  \node at (0.68,0.3) {$\frac{\pi}{3}$};

  % 点
  \fill (A) circle (1.5pt);

\end{tikzpicture}
    と表される. この点の原点からの距離は$2$,正の$x$軸から測った角度は$\pi/3$である.
  \subsection{極形式}
    一般に,原点からの距離を$\rho$,正の$x$軸から反時計回りに測った角度を$\theta$とすると,
    点の直交座標$(x,y)$は
    \[
      x=\rho\cos\theta,\qquad y=\rho\sin\theta
    \]
    と表される. ここで$\rho=\sqrt{x^2+y^2}$である.
    原点以外では$\rho>0$であり,$\theta$に$2\pi$の整数倍を加えても同じ点を表す.
    原点では$\rho=0$となり,角度は一意に定まらない.
